\chapter{The research behind Humanvoice}
\label{ch:literature-foundations}

An economist, referee, or executive does not read a technical document one
sentence at a time. The reader builds a working picture of the argument, checks
the evidence, and decides whether the requested action follows. Language-model
studies cover only part of that work. Humanvoice also draws on technical
communication, reading, feedback, human--computer interaction, and scholarly
integrity, fields that have studied those reader and writer decisions for much
longer.

Those fields do not give us direct efficacy evidence. Each contributes a way to
describe a reader, a revision, an interaction, or an integrity boundary. Before
we turn any lesson into a product claim, we need an observation that could test
it in technical documents.

The common idea is a loop rather than a score. Figure
\ref{fig:reader-author-loop} locates the judgments: the author owns meaning,
the reader owns use, and the system records the path between them.

\hvFigureReaderLoop

\section{Technical prose is a genre, not a neutral container}
\label{sec:lit-genre}

The first useful correction comes from register theory. Biber's analysis of
variation across speech and writing ties linguistic choices to situation,
audience, and purpose
\citep{biber1988variation}. A technical proposal is not a bad version of a
conversation, and a readable economics paper is not the same object as a
popular article. It uses equations, defined terms, cautious causal language,
and references because those forms carry specific obligations.

That observation changes the target for a writing tool. If a diagnostic rewards
short sentences without checking whether a definition has been separated from
its condition, it may move the text toward a different genre. If it removes
all repetition, it may erase the repeated noun that lets a reader track an
actor across a derivation. The right question is not ``does the document look
like ordinary speech?'' It is ``does the document give this reader the cues
needed for this technical task?''

Coh-Metrix makes the distinction measurable. Graesser, McNamara, Louwerse,
and Cai separate cohesion and language dimensions that syllable-based formulas
do not see \citep{graesser2004cohmetrix}. A pronoun with no clear antecedent,
a chain of nouns that changes referent, or a causal connector whose second
clause does not follow from the first can all increase reader work while
leaving sentence length unchanged. The CLEAR corpus extends the point with
human ease-of-reading judgments and shows that learned, multidimensional
features predict those judgments better than classical readability formulas
\citep{crossley2022clear}.

The evidence on scientific writing supplies a second warning. Scientific
abstracts have become harder to read over time as terminology and sentence
complexity rise \citep{plavensigray2017readability}. Specialized terminology
also carries a measurable communication cost: papers with more specialized
terms in titles and abstracts receive fewer citations in the study by
Mart\'{\i}nez and Mammola \citep{martinez2021terminology}. Neither result says
that technical terms should be removed. It says that a writer should pay for a
term with a definition, a reason for using it, or a later calculation that
needs it.

\begin{table}[H]
\centering
\small
\begin{tabularx}{\textwidth}{@{}p{0.23\textwidth}L p{0.27\textwidth}@{}}
\toprule
Reading property & What a genre-aware account notices & Humanvoice consequence \\
\midrule
Referential cohesion & Nouns, pronouns, and labels keep the same actor or
quantity in view & Detect a possible break, show the surrounding passage, and
ask the author to repair or defend it \\
Discourse connection & A result, reason, qualification, and implication have
an interpretable relation & Keep findings tied to the decision question, not
to an isolated sentence score \\
Terminology & Specialized words can be necessary but impose a knowledge cost &
Define a term at first use or show why the intended reader already knows it \\
Register & Formality and sentence shape vary by genre and audience & Compare
against an economics reference corpus, not a universal ``human'' baseline \\
\bottomrule
\end{tabularx}
\caption{The reading literature changes what a useful diagnostic is allowed to
mean.}
\label{tab:genre-consequences}
\end{table}

These are not four automatic scores. They are four reasons for locating a
passage and asking a substantive question. The author's answer may be that a
term is standard in the field, that a repeated label is needed for a theorem,
or that the paragraph belongs in a methods appendix. A diagnostic earns its
place by making that judgment easier, not by making the judgment disappear.

\section{Readers use position and effort as evidence}
\label{sec:lit-reader-expectations}

Consider two sentences with the same facts. The first ends with the method and
buries the result in its middle. The second starts from what the previous
sentence established and ends on the result the reader should carry forward.
No vocabulary list can explain why the second is easier to use. The difference
lies in where the sentence asks the reader to look for continuity and emphasis.

Gopen and Swan describe this problem through reader expectations: readers tend
to treat the opening of a sentence as contextual material and its closing
position as a place of emphasis; they also expect grammatical subjects and
verbs to reveal the story's actors and actions
\citep{gopen1990science}. Their account is an influential editorial model, not
a controlled demonstration that one arrangement improves every technical
document. It earns a place here because it yields specific, contestable
questions. What does this sentence connect to? Which actor performs the main
action? What new information receives emphasis? A Humanvoice finding can point
to a mismatch between those cues and the declared argument without pretending
to calculate comprehension.

The model also explains why noun-heavy prose feels bloodless. ``Validation of
the result was performed by the review layer'' makes the reader recover both
actor and action. ``The reviewer validated the result'' supplies them in the
grammar. Nominalization is not an error in itself: technical prose sometimes
needs a stable noun for a process that later becomes an object of analysis.
The diagnostic question is whether the noun helps the argument keep an object
in view or merely forces a weak verb to carry an action that could have been
stated directly.

Oppenheimer supplies adjacent experimental evidence for a related effect. In
several studies, replacing unnecessarily complex wording with simpler wording
changed judgments of the author; needless complexity did not make the author
appear more intelligent \citep{oppenheimer2006erudite}. The study concerns
processing fluency and judged intelligence, not acceptance of mathematical
monographs by expert executives. It therefore supports a restrained inference:
complex vocabulary should earn its cost. It does not support deleting domain
terms, shortening every sentence, or treating a simple-word ratio as a quality
score.

Together, these sources sharpen the product design. A finding should show the
passage and the relation it may obscure: actor to action, old information to
new, evidence to conclusion, or term to purpose. The interface should offer a
question such as ``what should the reader carry from the end of this sentence
into the next one?'' rather than an instruction such as ``simplify.'' The
reader study can then test whether the revised passage improves reconstruction
or reduces time. That later observation, not the editorial maxim, determines
whether the finding belongs in the product.

\section{Useful feedback changes the next draft}
\label{sec:lit-feedback-foundations}

An author can receive a beautifully specific comment and still write the next
draft unchanged. Feedback research therefore supports the product's revision
loop only conditionally: a comment is useful when it changes a later choice in
a way that helps the declared task. Specificity and politeness alone are not
evidence.

Nair and colleagues evaluate generated feedback on an economic essay task by
looking at simulated revisions rather than by asking whether the comments
sound detailed \citep{nair2024closing}. Rashkin and colleagues construct
stories with deliberately corrupted passages and show that models can find
local defects while missing the most important problem or selecting praise
where criticism is needed \citep{rashkin2025story}. Li and colleagues find
that a detailed rubric can make a model a useful revision assessor, but that
agreement varies with writer proficiency \citep{li2024revisions}. Behzad and
colleagues report that model and human feedback can complement one another,
while personalized feedback remains difficult \citep{behzad2024feedback}.

Taken together, these studies imply three distinct units of analysis:

\begin{enumerate}
\item the \emph{finding}, which says what passage deserves attention;
\item the \emph{revision}, which records what the author changed and why; and
\item the \emph{outcome}, which says whether the intended reader could complete
      the task with less burden and no meaning error.
\end{enumerate}

Keep the three observations separate. A high-precision finding can be useless
if it concerns a low-consequence phrase. A beautiful revision can be harmful if
it changes the denominator in a table. A reader can accept a document for the
wrong reason. The evaluation chapter therefore treats finding, revision, and
reader outcome as a chain with separate records and separate uncertainty.

For document $d$, revision path $h$, and reader task $q$, write
\begin{equation}
  Y(d,h,q) = \bigl(F(d),\; C(h\mid F),\; A(q\mid d,h)\bigr),
  \label{eq:lit-three-outcomes}
\end{equation}
where $F$ is the set of located findings, $C$ is the author's change and
explanation, and $A$ is the reader's ability to act. Equation
\eqref{eq:lit-three-outcomes} is a bookkeeping identity for the product, not a
psychological model. It prevents a benchmark of $F$ from being reported as a
benchmark of $A$.

The edit study by Chakrabarty, Laban, and Wu is useful because it measures the
middle link directly. Professional writers labeled span-level idiosyncrasies,
an automated detector recovered only part of those spans, and edited text
outperformed raw model output while remaining behind professional writing
\citep{chakrabarty2024salvage}. The result supports an edit workflow, not a
one-shot humanizer. It also gives a realistic expectation for the first
release: finding precision will be imperfect, so the interface must expose
context and allow rejection without penalty.

\section{The reader is part of the intervention}
\label{sec:lit-reader-interaction}

Human--AI interaction research prevents a common experimental mistake: treating
the final text as if it were produced by a single prompt. CoAuthor records
keystrokes and suggestion use in a collaborative writing setting
\citep{lee2022coauthor}. Wordcraft shows how writers use a model to explore,
select, and revise alternatives rather than simply accept its first answer
\citep{yuan2022wordcraft}. Mysore and colleagues' in-the-wild traces identify
recurring paths in which users revise their intent, ask questions, adjust
style, and inject content \citep{mysore2025paths}. A document-level experiment
that records only the before and after versions loses the process variables
that explain both benefit and harm.

When a reader asks a question, the author should be able to see which suggestion
was offered, which one was accepted or rejected, which protected object was
touched, and whether the question caused a second revision. Those records
support analysis; they should not become visible workflow vocabulary in the
finished proposal. The reader should see a clear paragraph, not a history of
prompts.

Agency has a second dimension. The author can delegate wording while retaining
responsibility for a claim, or can delegate a structural diagnosis while
retaining the final edit. The system should ask which delegation occurred
because disclosure and accountability depend on it. It should not infer
authorship from the text. The detector studies in Chapter~\ref{sec:literature}
show why: formal or non-native writers can be flagged as machine-generated,
while paraphrase can evade several detector families
\citep{liang2023detectors,krishna2023paraphrasing}.

\section{Integrity is a property of the path}
\label{sec:lit-integrity}

Citation, disclosure, and provenance are often discussed as policy questions,
but each has a concrete technical counterpart. A citation can have a valid
identifier and still fail to support the sentence. A disclosure can be present
and still leave the reader unable to tell which passages were assisted. A
version history can show that text changed and still omit the reason for a
changed number.

Mugaanyi and colleagues' evaluation of references generated by ChatGPT finds
unverifiable and mismatched citations often enough to make mechanical
resolution and human checking necessary \citep{mugaanyi2024citation}. ICMJE's
recommendations place accuracy, integrity, originality, and responsibility
with human authors, while treating a language model as a tool rather than an
author \citep{icmje2025ai}. NIST's AI Risk Management Framework gives a broader
vocabulary for validity, reliability, privacy, transparency, and accountability
\citep{nist2023airmf}.

Those sources do not prescribe Humanvoice's interface. They do establish a
design relation:

\begin{equation}
  \begin{aligned}
  \text{responsible revision} ={}&\text{located source}
      +\text{declared assistance}\\[-1pt]
      &+\text{human ownership of the decision}.
  \end{aligned}
  \label{eq:lit-integrity-relation}
\end{equation}

The relation is not a score and the plus signs are not arithmetic. Each term is
necessary for a reader to reconstruct what happened. If the source cannot be
located, a changed claim cannot be examined. If assistance is undisclosed,
the reader cannot interpret the interaction or the responsibility boundary. If
no human owns the decision, a technically fluent output has no accountable
author.

\section{Evaluation traditions and their blind spots}
\label{sec:lit-evaluation-traditions}

The literature uses several familiar endpoints: time, rubric scores, pairwise
preference, agreement with experts, and task completion. They answer different
questions. Noy and Zhang report faster and better-scored occupation-specific
writing with ChatGPT \citep{noy2023experimental}; Brynjolfsson, Li, and
Raymond find productivity gains in customer support with heterogeneous effects
by worker experience \citep{brynjolfsson2025work}; and Dell'Acqua and
colleagues show that assistance can help inside a capability frontier and hurt
outside it \citep{dellacqua2026jagged}. These studies are valuable evidence
about task-level assistance. They are not estimates of reader acceptance for a
technical proposal.

The LLM-as-judge literature supplies a parallel warning. Zheng and colleagues
measure position and verbosity bias in model evaluation
\citep{zheng2023judging}. Panickssery, Bowman, and Feng show that models can
recognize and favor their own generations \citep{panickssery2024llm}. Dubois
and colleagues show that length control changes automatic rankings and improves
agreement with human judgments \citep{dubois2024length}. Chakrabarty and
colleagues find that model judges are generous toward generated creative text
where expert writers are stricter \citep{chakrabarty2024art}.

The evaluation design follows from the blind spots rather than from a desire
to collect every metric. The primary endpoint is whether a named reader can
make the declared decision with less total attention and no protected-object
error. Time and feedback rounds explain cost. Pairwise model judgments help
prioritize repairs after calibration. Surface diagnostics describe style drift.
None of the secondary measures can replace the primary endpoint.

\section{Foundations of writing as planned work}
\label{sec:lit-writing-foundations}

The product claim also rests on an older literature that is easy to lose when
the discussion is narrowed to language models. Flower and Hayes describe
writing as a recursive process of planning, translating, and reviewing under a
writer's goals and a representation of the rhetorical problem
\citep{flower1981cognitive}. The useful unit is therefore not a sentence in
isolation. It is a choice made while the writer is trying to satisfy a reader,
an evidence standard, and a purpose at the same time. A tool that improves a
sentence while hiding which goal it serves may reduce local friction while
increasing the next planning cost.

Bereiter and Scardamalia's distinction between knowledge telling and knowledge
transforming makes the same point in a different vocabulary
\citep{bereiter1987psychology}. Knowledge telling retrieves propositions and
places them on the page. Knowledge transforming changes the relation among
claims, reasons, evidence, and the problem being solved. Expert technical
documents require the second operation: a result must be selected, situated,
qualified, and connected to an action. A fluent assistant can make knowledge
telling faster without making the transformation happen. Humanvoice should
therefore expose a missing relation as a question, not fill it with plausible
text.

Genre analysis supplies the social side of the model. Swales treats research
writing as a sequence of moves that establishes a territory, identifies a
gap, occupies that gap, and qualifies the resulting claim
\citep{swales1990genre}. The exact moves vary by discipline and document type,
but the implication is stable: a proposal, methods note, and executive
memorandum have different obligations even when they share vocabulary. A
universal ``human'' baseline is not a defensible target. The reference corpus
and reading brief must name the genre, the decision, and the audience whose
expectations make a move useful.

Hyland's account of metadiscourse adds a reader-facing layer
\citep{hyland2005metadiscourse}. Writers use signposts, stance, and
engagement to tell readers how to interpret a claim, how strongly to take it,
and where to look next. Some of the patterns that a generic humanizer marks as
repetition are doing this work. Removing every ``however,'' condition, or
roadmap sentence can make a document shorter while making its inferential
path harder to follow. A diagnostic must ask whether the device is redundant
in this passage, not whether the device is common in model output.

Feedback research then explains why the proposed loop ends with a reader. Hattie
and Timperley distinguish feedback about the task, the process, and the
self-regulatory activity around the work; the effects depend on timing,
specificity, and how the recipient uses the information
\citep{hattie2007feedback}. Nicol and Macfarlane-Dick emphasize that good
feedback helps a learner generate and use judgments rather than merely receive
a correction \citep{nicol2006formative}. For a technical author, the analogue
is an editorial question that can be investigated and accepted or rejected,
not a replacement paragraph that silently transfers responsibility to a
system.

\begin{table}[H]
\centering
\small
\begin{tabularx}{\textwidth}{@{}p{0.22\textwidth}L p{0.27\textwidth}@{}}
\toprule
Tradition & Observation relevant to the product & Boundary on the claim \\
\midrule
Writing as a recursive process & The author plans for a rhetorical problem,
translates ideas, and reviews against goals. & A sentence score cannot stand in
for a model of the document's purpose. \\
Knowledge transforming & Expertise changes relations among claims, reasons,
evidence, and consequences. & Faster drafting does not show that the relation
was transformed correctly. \\
Genre and metadiscourse & Research and decision genres use recognizable moves,
signposts, stance, and engagement. & A universal style baseline can erase a
necessary disciplinary cue. \\
Formative feedback & Feedback is useful when it helps the writer generate and
use a better judgment on the next attempt. & A detailed comment is not evidence
of uptake, transfer, or reader acceptance. \\
\bottomrule
\end{tabularx}
\caption{Foundational writing research supplies mechanisms and boundaries for
the product hypothesis.}
\label{tab:writing-foundations}
\end{table}

Consider a small worked contrast. A local checker sees the sentence ``The
model performs well'' and can suggest a more active verb. The writing-process
literature asks a prior question: what is the rhetorical problem? If the
reader must decide whether to fund a larger run, the sentence needs a declared
comparison, an outcome, and a condition. A defensible revision might be
``On the held-out documents, the constrained reconstruction reduces median
reader time by 11 minutes, with no protected-object changes.'' That revision
is not justified by fluency. It is justified only if the comparison, unit,
sample, and qualification exist in the source. The same sentence would be
wrong in a methods section if the 11 minutes came from an illustrative fixture
rather than observed readers.

The tempting inference is that a model which proposes the second sentence has
understood the argument. It may instead have supplied a familiar empirical
shape. Humanvoice makes the distinction operational: the finding points to the
span and the missing question; the author supplies or rejects the evidence; the
protected comparison records the change; and the reader instrument tests
whether the resulting claim can be used. The literature makes this division of
labor plausible. It does not make the proposed revision true.

These foundations change the implementation in four ways. First, the source
snapshot stores the reading brief and declared decision alongside text, so a
finding can be interpreted as a rhetorical problem rather than a lexical
anomaly. Second, the diagnostic vocabulary includes relations among actors,
evidence, qualifications, and actions, while retaining simple surface rules
as low-cost candidates. Third, the author log records why a suggestion was
accepted, rejected, or deferred; this is the observable trace of feedback use,
not a productivity badge. Fourth, the reader exercise is part of the workflow,
because transfer to a later decision is the outcome that matters.

The foundations also limit the first release. Humanvoice should not attempt to
infer a complete rhetorical move sequence for every discipline, score a
writer's expertise, or prescribe one voice. It can identify a possible missing
actor, relation, or qualification and show the evidence behind that question.
It can compare protected objects and ask a reader what remains unclear. The
author and the domain reader still supply the interpretation. This narrow
boundary is what makes a feasibility study possible without pretending to
automate composition theory.

\section{Gaps in the review}
\label{sec:lit-gaps}

A skeptical panel should ask what is missing. The current snapshot has live
broad and challenge searches, public imports, a provenance-complete bounded
RePEc/IDEAS specialist slice, and resolved Crossref identities for every
DOI-bearing bibliography row. Broad retrieval passes, but four of eighteen
known items still evade the deliberately indirect challenge queries. The
snapshot also has no independent screening of all load-bearing rows and no
full-text verification or design-specific appraisal for every claim anchor.
Those remaining gaps are explicit, not hidden assumptions.

The topic also has adjacent literatures that deserve deliberate treatment in
the next search pass: technical-communication studies of expert review,
composition research on revision and feedback uptake, information-retrieval
work on citation entailment, human-factors work on decision aids, and empirical
studies of privacy-preserving local inference. The product can proceed to a
bounded feasibility build while those searches continue because the current
request is for implementation learning. It cannot present the resulting
pilot as a completed systematic review or as proof that the workflow improves
all expert writing.

The omission register in Appendix~\ref{app:literature-matrix} records the
question, likely reviewer objection, source status, and next action for each
high-risk omission. This is a better scholarly posture than filling the
bibliography with names that have not been read. The next chapter applies the
same standard to software: an available package earns a place in the product
only after its capability, license, failure behavior, and evidence burden are
visible.
